% arXiv preprint skeleton -- figures and tables are produced by `python cli.py paper-figs`
% (copy results/benchmark/figures/* next to this file, or adjust \graphicspath).
\documentclass[11pt]{article}
\usepackage[margin=1in]{geometry}
\usepackage{amsmath,amssymb,graphicx,booktabs,hyperref,siunitx,xcolor}
\graphicspath{{../../results/benchmark/figures/}}
\newcommand{\todo}[1]{\textcolor{red}{[TODO: #1]}}
% Generated by scripts/make_numbers.py from results/benchmark/*.csv -- do not edit by hand.
\newcommand{\RzeroKL}{99.1}  % R0 membrane+bending, degree 4: non-affine energy captured [%]
\newcommand{\RzeroKLDims}{23}  % R0 membrane+bending degree 4: dimensions
\newcommand{\RzeroGeneric}{96.2}  % R0 generic polynomial, degree 4 [%]
\newcommand{\RzeroGenericDims}{24}  % R0 generic polynomial degree 4: dimensions
\newcommand{\RzeroMembrane}{45.6}  % R0 membrane only, degree 4 [%]
\newcommand{\RzeroKLResidNm}{0.58}  % R0 membrane+bending: median RMS left [nm]
\newcommand{\RzeroGenericResidNm}{1.18}  % R0 generic polynomial: median RMS left [nm]
\newcommand{\NdmaxOurs}{64}  % E1: FEM runs to reach median R2(d_max) >= 0.95, ours
\newcommand{\NdmaxBaseline}{512}  % E1: same, best learned baseline
\newcommand{\EoneOursThirtyTwo}{0.925}  % E1: R2(d_max) at N=32, ours
\newcommand{\EoneBaseThirtyTwo}{0.791}  % E1: R2(d_max) at N=32, best learned baseline
\newcommand{\EoneMembraneDmaxSixtyFour}{0.959}  % E1: R2(d_max) at N=64, membrane-only basis
\newcommand{\EoneOursDmaxSixtyFour}{0.956}  % E1: R2(d_max) at N=64, ours
\newcommand{\EoneNoPriorDmaxSixtyFour}{0.887}  % E1: R2(d_max) at N=64, ours without analytic prior
\newcommand{\EoneOursResidOneTwoEight}{0.838}  % E1: R2(resid_3sigma) at N=128, ours
\newcommand{\EoneMembraneResidOneTwoEight}{0.670}  % E1: R2(resid_3sigma) at N=128, membrane-only
\newcommand{\EoneOursFieldFiveTwelve}{15.5}  % E1: field RMSE at N=512, ours [nm]
\newcommand{\EoneMembraneFieldFiveTwelve}{18.5}  % E1: field RMSE at N=512, membrane-only [nm]
\newcommand{\EtwoOursFour}{2.5}  % E2: 4 marks, 2 nm noise, 256-run prior: median RMSE [nm]
\newcommand{\EtwoOursFourSmallPrior}{4.4}  % E2: same with 64-run prior [nm]
\newcommand{\EtwoAffineFour}{8.1}  % E2: 4 marks, 2 nm noise, affine overlay model [nm]
\newcommand{\EtwoPriorOnly}{8.1}  % E2: 256-run prior, no marks [nm]
\newcommand{\EtwoPriorAffineFour}{2.6}  % E2: 4 marks, 2 nm noise, 256-run prior + affine fit [nm]
\newcommand{\EtwoOursNine}{2.5}  % E2: 9 marks, 2 nm noise, 256-run prior, ours [nm]
\newcommand{\EtwoPriorAffineNine}{2.4}  % E2: 9 marks, 2 nm noise, 256-run prior + affine fit [nm]
\newcommand{\EtwoOursSixteen}{1.4}  % E2: 16 marks, 2 nm noise, 256-run prior, ours [nm]
\newcommand{\EtwoPriorAffineSixteen}{2.2}  % E2: 16 marks, 2 nm noise, 256-run prior + affine fit [nm]
\newcommand{\EtwoPriorAffineFourSmallPrior}{4.4}  % E2: 4 marks, 2 nm noise, 64-run prior + affine fit [nm]
\newcommand{\EtwoOursNineSmallPrior}{3.6}  % E2: 9 marks, 2 nm noise, 64-run prior, ours [nm]
\newcommand{\EtwoPriorAffineNineSmallPrior}{4.4}  % E2: 9 marks, 2 nm noise, 64-run prior + affine fit [nm]
\newcommand{\EtwoOursSixteenSmallPrior}{1.7}  % E2: 16 marks, 2 nm noise, 64-run prior, ours [nm]
\newcommand{\EtwoPriorAffineSixteenSmallPrior}{3.9}  % E2: 16 marks, 2 nm noise, 64-run prior + affine fit [nm]
\newcommand{\EtwoPolyTwoNine}{7.0}  % E2: 9 marks, 2 nm noise, 2nd-order correction [nm]
\newcommand{\EtwoPolyThreeSixteen}{2.2}  % E2: 16 marks, 2 nm noise, 3rd-order correction [nm]
\newcommand{\EtwoAffineSixteen}{5.8}  % E2: 16 marks, 2 nm noise, affine [nm]
\newcommand{\EthreeOodDTOurs}{0.96}  % E3 ood_dT_head_gt40: R2(d_max), ours
\newcommand{\EthreeOodDTBaseMin}{0.69}  % E3 ood_dT_head_gt40: R2(d_max), worst learned baseline
\newcommand{\EthreeOodDTBaseMax}{0.82}  % E3 ood_dT_head_gt40: R2(d_max), best learned baseline
\newcommand{\EthreeOodLOurs}{0.96}  % E3 ood_L_die_gt14: R2(d_max), ours
\newcommand{\EthreeOodLBaseMin}{0.66}  % E3 ood_L_die_gt14: R2(d_max), worst learned baseline
\newcommand{\EthreeOodLBaseMax}{0.92}  % E3 ood_L_die_gt14: R2(d_max), best learned baseline
\newcommand{\EthreeNoiseFiveOurs}{0.99}  % E3 train_noise_5nm: R2(d_max), ours
\newcommand{\EthreeNoiseFiveBaseMin}{0.83}  % E3 train_noise_5nm: R2(d_max), worst learned baseline
\newcommand{\EthreeNoiseFiveBaseMax}{0.90}  % E3 train_noise_5nm: R2(d_max), best learned baseline
\newcommand{\PinnFieldPhys}{92}  % E4: PINN field RMSE at N=64, with physics loss [nm]
\newcommand{\PinnFieldNoPhys}{88}  % E4: PINN field RMSE at N=64, physics loss removed [nm]
\newcommand{\PinnDmaxSixtyFour}{0.83}  % E4: PINN R2(d_max) at N=64
\newcommand{\NParams}{16}  % swept process/geometry/material parameters
\newcommand{\NTrainRuns}{1024}  % FEM training runs (Sobol)
\newcommand{\NHoldoutRuns}{128}  % FEM holdout runs (Latin hypercube)
\newcommand{\NPads}{21}  % pads per die edge (n x n grid)
\newcommand{\MeshCoarse}{12}  % elements along the die edge in the datasets
\newcommand{\RunTimeMedianS}{23}  % median wall time per FEM run [s]
\newcommand{\MeshFine}{28}  % finest mesh in the convergence study
\newcommand{\ConvCoarsePct}{0.2}  % worst m_x/d_max change, coarse vs finest mesh [%]
\newcommand{\DofsCoarse}{99,858}  % displacement DOFs, die + wafer, coarse mesh
\newcommand{\BondLayerEffectPct}{0.3}  % largest d_max change when the bond layers get Si properties [%]
\newcommand{\SiE}{130}  % Si Young's modulus [GPa]
\newcommand{\SiNu}{0.28}  % Si Poisson ratio
\newcommand{\SiAlpha}{2.6}  % Si thermal expansion [ppm/K]
\newcommand{\CuE}{120}  % Cu Young's modulus [GPa]
\newcommand{\CuAlpha}{16.7}  % Cu thermal expansion [ppm/K]
\newcommand{\Tref}{20}  % reference (pad patterning) temperature [C]
\newcommand{\WaferPatchFactor}{2}  % wafer patch edge / die edge
\newcommand{\NzDie}{4}  % elements through the die thickness
\newcommand{\NzWafer}{4}  % elements through the wafer thickness
\newcommand{\NRing}{6}  % elements across the wafer margin
\newcommand{\NTimeSteps}{20}  % thermal time steps per run
\newcommand{\TFirstUs}{1}  % first thermal time step [us]
\newcommand{\VerifTolPct}{1}  % tolerance of the analytic checks in tests/test_verification.py [%]
  % generated by scripts/make_numbers.py -- every number in the text is a macro from here

\title{Physics-Structured Learning of Bonding-Induced Overlay Distortion\\ in Data-Scarce Die-to-Wafer Hybrid Bonding}
\author{Perry (Chao-Pei) Lu}
\date{}

\begin{document}
\maketitle

\begin{abstract}
Die-to-wafer (D2W) hybrid bonding at sub-\SI{2}{\micro\meter} pitch leaves an overlay budget of a few hundred
nanometers, yet the data needed to predict die shift are scarce on both ends: each high-fidelity
thermo-mechanical simulation is expensive in a 16-dimensional process space, and production metrology
measures only a handful of overlay marks per die. We propose a physics-informed learning framework that
encodes the kinematics of the bonded die surface---the response of a free plane-stress membrane to an
isotropic eigenstrain plus a curl-free bending field (Kirchhoff--Love)---together with a closed-form
magnification prior. Against a 3D transient finite-element model with bond-wave locking, this reduced physics
represents 99.1\% of the non-affine overlay energy with 23 degrees of freedom, versus 96.2\% for a generic
polynomial with 24, and it shows that the non-uniform die shift is bending-dominated (membrane alone: 45.6\%).
A hard-constrained realization (physics basis with Gaussian-process coefficients) reaches a holdout
$R^2\geq0.95$ for the maximum pad misalignment with 64 FEM runs, where the best hybrid baseline (analytic model
plus GP residual) needs 512, and it keeps $R^2=0.96$ when extrapolating to hotter bond heads, against 0.82--0.83
for the baselines. For sparse metrology, the physics prior lowers the full-die reconstruction error of the industry six-parameter
overlay model by a factor of \EtwoPriorGainFourSmall--\EtwoPriorGainFourLarge\ with \EtwoKFour~marks, and a Bayesian
reconstruction in the physics span lowers it by a further \EtwoBasisGainSixteenLarge--\EtwoBasisGainSixteenSmall\ when
\EtwoKSixteen~marks are available. A soft-constrained parametric PINN with
the same physics is competitive only at the smallest data size, and its physics residual adds no measurable
accuracy---we analyze why. Code, data and the pinned FEM container are released.
\end{abstract}

\section{Introduction}
\begin{itemize}
  \item D2W hybrid bonding overlay budget ($<\SI{250}{\nano\meter}$ at \SI{2}{\micro\meter} pitch \cite{bex2024}); die shift from
        bond-head heating, die bow and head compliance; translation/rotation corrected by the bonder, the rest is not.
  \item Data scarcity on two levels: (i) simulation cost in a 16-D process space; (ii) sparse overlay marks in production.
  \item Existing practice: 4--6 parameter linear overlay models \cite{bex2024}; data-driven ML on overlay \cite{james2024};
        PINNs in packaging have targeted joint strength \cite{mei2026sinteredag} and warpage \cite{warpagepinn2026}.
  \item Contributions: (1) reduced physics of the bonded die surface and its validation against 3D transient FEM (R0);
        (2) two physics-informed realizations (hard basis + GP; parametric PINN); (3) data-efficiency benchmark against
        fair hybrid baselines (E1, E4); (4) sparse-mark reconstruction that beats the industry overlay model (E2);
        (5) robustness to noise and extrapolation (E3); (6) open, reproducible pipeline.
\end{itemize}

\section{Related work}
\todo{D2W overlay error budgets \cite{bex2024}; hybrid-bonding yield models \cite{chen2025yapplus}; die shift in fan-out
packaging \cite{sung2026dieshift}; physics/data hybrid models and digital twins for 3D IC packaging \cite{datta2026digitaltwin};
W2W bonding mechanics \cite{ghosh2026w2w}; ML for bonding overlay \cite{james2024}; PINNs \cite{raissi2019pinn} and PINN
applications in packaging and contact mechanics \cite{warpagepinn2026,mei2026sinteredag,xia2024pinncontact}.}

\section{High-fidelity reference model}\label{sec:reference}
All learning and all evaluation in this paper use a three-dimensional transient finite-element model of the bonding
moment as the reference. It is more complete than the reduced physics of Section~\ref{sec:method}: it resolves the
temperature field in time, the die and the wafer as separate elastic bodies, and the progressive locking of pads by the
bond wave. It has not been compared with measured overlay; every accuracy figure below is relative to this model
(Section~\ref{sec:discussion}).

\paragraph{Geometry and materials.} A square silicon die of edge $L$ and thickness $t$ faces a square patch of the target
wafer, \WaferPatchFactor~times wider than the die, of thickness $t_w$. Each bonding surface carries a bond layer of
thickness $t_b$, a mixture of dielectric and Cu pads homogenized by the rule of mixtures (stiffness-weighted for the
thermal expansion) with Cu area fraction $\phi_{\mathrm{Cu}}$. Silicon is isotropic with $E=\SI{\SiE}{\giga\pascal}$,
$\nu=\SiNu$ and $\alpha_{\mathrm{Si}}=\SI{\SiAlpha}{ppm/K}$; Cu has $E=\SI{\CuE}{\giga\pascal}$ and
$\alpha=\SI{\CuAlpha}{ppm/K}$. Pad patterns are aligned and stress-free at $T_{\mathrm{ref}}=\SI{\Tref}{\celsius}$.

\paragraph{Thermal model.} At contact the die is at the bond-head temperature
$T_{\mathrm{head}}=T_{\mathrm{chuck}}+\Delta T_{\mathrm{head}}$ and the wafer at the chuck temperature. Heat conduction
through die, bond layers and wafer (perfect contact at the interface) is integrated with backward Euler over the time the
bond wave needs to reach the farthest die corner; the die top is held at $T_{\mathrm{head}}$ and the wafer bottom at
$T_{\mathrm{chuck}}$. The \NTimeSteps~time steps grow geometrically from \SI{\TFirstUs}{\micro\second}, because the thermal
transient after contact is much shorter than the passage of the wave.

\paragraph{Mechanical model.} Before a pad bonds, die and wafer are not connected, so each is solved as its own linear
thermo-elastic body under the instantaneous temperature field (quasi-static). The eigenstrain is
$\alpha\,(T-T_{\mathrm{ref}})$; in the die it includes a through-thickness gradient that gives the die its initial
curvature $\kappa$. The bond head acts on the die top as a normal spring of stiffness $k_{\mathrm{head}}$, between a free
die ($k_{\mathrm{head}}\to0$) and a rigid flat head that flattens the die ($k_{\mathrm{head}}\to\infty$). The bond force
$F$ enters as an equal and opposite pressure $F/L^2$ on the die top and bottom, and on the wafer within the die
footprint. The wafer bottom follows the thermal expansion of the chuck, $\alpha_c(T_{\mathrm{chuck}}-T_{\mathrm{ref}})$,
and the patch edges follow the expansion of the surrounding wafer. The in-plane rigid motion of the die is removed by
point constraints at the centre of its bonding surface; this is the anchoring convention discussed in
Section~\ref{sec:quantity}.

\paragraph{Bond wave and overlay.} The bond front starts at $(x_0,y_0)L$ and spreads radially at speed $v$. A pad at
distance $r$ from the origin locks at $t=r/v$, and its overlay is the in-plane displacement of the die pad minus that of
the wafer pad at that time, interpolated between time steps. The rigid translation and rotation present before contact
are subtracted, as the bonder's alignment would remove them. The overlay is recorded on an $\NPads\times\NPads$ pad grid
covering the die.

\paragraph{Discretization and verification.} The mesh is structured and hexahedral, with quadratic displacement
elements: \MeshCoarse~elements along the die edge in the datasets, \NRing~across the wafer margin, \NzDie~and
\NzWafer~through the die and wafer thicknesses and one per bond layer, about \DofsCoarse~displacement unknowns in total.
Quadratic elements are needed because linear hexahedra lock in bending at these aspect ratios. A mesh study at the
default parameters and at the thin-layer corner of the parameter ranges (thinnest bond layer and die, largest die)
shows that this mesh changes $m_x$ and $d_{\max}$ by at most \ConvCoarsePct\% relative to
$\MeshFine$ elements along the edge. Automated tests check the model against closed-form limits within
\VerifTolPct\%: a uniform thermal mismatch gives the exact magnification, a free curved die gives the plate bending
strain, and a rigid head suppresses it; zero load gives zero overlay, and a mirrored wave origin mirrors the translation.

\paragraph{Parameters and datasets.} The \NParams~swept parameters and their ranges are listed in
Table~\ref{tab:params}. Replacing the bond-layer properties by those of silicon changes $d_{\max}$ by at most
\BondLayerEffectPct\%, so the four bond-layer parameters are expected to matter little; they are kept to test that the
learned models find this. The training set has \NTrainRuns~runs on a scrambled Sobol sequence and the holdout set
\NHoldoutRuns~runs on an independent Latin hypercube; holdout runs are never used for training or model selection. A
run takes about \RunTimeMedianS~s (median, two threads).

\begin{table}[t]\centering\small
\caption{Swept parameters of the reference model. Log-scaled parameters are sampled uniformly in $\log$.}\label{tab:params}
\input{../../results/benchmark/figures/table_params.tex}
\end{table}

\section{Method}\label{sec:method}

\subsection{Predicted quantity}\label{sec:quantity}
For one parameter set $\mathbf{p}\in\mathbb{R}^{16}$ the reference model returns the overlay vector
$\mathbf{d}_i\in\mathbb{R}^2$ at every pad $i$ of an $n\times n$ grid on the bonding interface: the in-plane
displacement of the die pad minus that of the facing wafer pad, each taken at the moment the bond wave locks the
pad. Two conventions of the reference model fix what $\mathbf{d}$ contains. First, the translation and rotation present
before contact are removed, as the bonder's alignment would. Second, the die's in-plane rigid motion is anchored at
its centre. Translation and rotation therefore stay close to zero by construction, and $\mathbf{d}$ is the
\emph{bonding-induced overlay distortion}: magnification plus a non-uniform remainder. Placement error of the bonder
is not modelled; it can be superposed afterwards. We do not call $\mathbf{d}$ ``die shift'' for this reason.

All field models work on the reference square. With $L$ the die edge length, a pad at $\mathbf{x}$ maps to
$\boldsymbol\xi=\mathbf{x}/(L/2)\in[-1,1]^2$, and the field is stored in strain units,
\begin{equation}
  \mathbf{f}(\boldsymbol\xi;\mathbf{p}) = \frac{\mathbf{d}}{L/2},
\end{equation}
so that \SI{1}{ppm} of uniform strain moves a pad at the die edge by $L/2$~nm when $L/2$ is in millimetres. In these
units fields of dies of different size are directly comparable. Scalar metrics are derived from a predicted field with
the same least-squares overlay fit that is applied to the reference data,
\begin{equation}\label{eq:overlayfit}
  \mathbf{d}_i = \mathbf{t} + \theta\begin{pmatrix}-y_i\\ x_i\end{pmatrix}
               + \begin{pmatrix}m_x x_i\\ m_y y_i\end{pmatrix} + \mathbf{r}_i ,
\end{equation}
which gives the magnifications $m_x, m_y$, the non-uniform component $\mathrm{resid}_{3\sigma}=3\,\mathrm{std}(\mathbf{r})$,
and the largest pad misalignment $d_{\max}=\max_i\lVert\mathbf{d}_i\rVert$.

\subsection{Reduced physics of the bonding surface}\label{sec:reduced}
The die is thin compared with its edge length, so we describe its bonding surface with Kirchhoff--Love kinematics. The
in-plane displacement of the surface is a membrane part plus a bending part,
\begin{equation}\label{eq:kl}
  \mathbf{f} = \mathbf{u}_m + \nabla\phi, \qquad \phi = -\frac{t}{2}\,w ,
\end{equation}
where $t$ is the die thickness and $w$ the deflection of its mid-plane. The bending part is the gradient of a scalar
potential and therefore curl-free; no assumption on the bending loads (bond-head compliance, bow, through-thickness
temperature gradient) is needed for this structure.

The membrane part obeys plane-stress equilibrium of a free die loaded by an in-plane eigenstrain $e(\boldsymbol\xi)$,
written in stresses scaled by $E/(1-\nu^2)$:
\begin{equation}\label{eq:membrane}
  \nabla\cdot\boldsymbol\sigma = \mathbf{0}\ \text{in}\ \Omega,\qquad
  \boldsymbol\sigma\,\mathbf{n} = \mathbf{0}\ \text{on}\ \partial\Omega,\qquad
  \boldsymbol\sigma = (1-\nu)\,\boldsymbol\varepsilon(\mathbf{u}_m)
     + \nu\,\mathrm{tr}\,\boldsymbol\varepsilon(\mathbf{u}_m)\,\mathbf{I} - (1+\nu)\,e\,\mathbf{I},
\end{equation}
with $\Omega=[-1,1]^2$ and $\boldsymbol\varepsilon$ the symmetric gradient. We take the eigenstrain isotropic in the
plane: thermal mismatch, the flattening of a bowed die and the Poisson effect of the bond pressure all act equally in
both in-plane directions for an isotropic die. A uniform eigenstrain $e_0$ produces a stress-free state with
$\mathbf{u}_m=e_0\,\boldsymbol\xi$, i.e.\ pure isotropic magnification.

For that uniform part a closed-form estimate is available from the process parameters,
\begin{equation}\label{eq:prior}
  m_a(\mathbf{p}) = \alpha_{\mathrm{Si}}\,(T_{\mathrm{head}}-T_{\mathrm{ref}})
                  - \alpha_{c}\,(T_{\mathrm{chuck}}-T_{\mathrm{ref}})
                  + \kappa\,\frac{t+t_{b}}{2},
\end{equation}
the expansion of a die at the head temperature, minus the expansion the chuck imposes on the wafer, plus the surface
strain released when a die of curvature $\kappa$ is flattened ($t_b$: bond-layer thickness, $T_{\mathrm{ref}}$: pad
patterning temperature). We use $m_a$ as the prior on the uniform eigenstrain in both realizations below.

Equations~\eqref{eq:kl}--\eqref{eq:membrane} describe the \emph{shape} of admissible fields; they do not determine
$e(\boldsymbol\xi)$ or $w$ for a given $\mathbf{p}$, which depend on the transient temperature field and on contact.
That dependence is learned from the reference runs. The representation test R0 (Section~\ref{sec:experiments}) measures how much of the reference overlay
the admissible fields can represent.

\subsection{Physics-structured realization: basis and Gaussian processes}\label{sec:basis}
\paragraph{Basis.} The admissible fields of Section~\ref{sec:reduced} are spanned by three groups of columns evaluated at
the pads:
\begin{itemize}
  \item membrane responses $\mathbf{U}_{ij}$ to the eigenstrain monomials $e=\xi^i\eta^j$, $i+j\le4$, obtained by
        solving~\eqref{eq:membrane} once on the reference square with biquadratic finite elements (scikit-fem);
        the uniform response is set to its exact value $\mathbf{U}_{00}=\boldsymbol\xi$;
  \item bending fields $\nabla(\xi^i\eta^j)$, $2\le i+j\le4$;
  \item the three rigid modes (two translations and a rotation).
\end{itemize}
The membrane and bending columns other than $\mathbf{U}_{00}$ are projected out of the span of $\mathbf{U}_{00}$ and the
rigid modes and orthonormalized by a singular value decomposition, dropping directions whose singular value is below
$10^{-8}$ of the largest. The result is a matrix $\mathbf{B}=[\mathbf{U}_{00},\mathbf{Q},\mathbf{R}]$ in which the first
coefficient of any field is its isotropic magnification. The orthonormalization matters: with nearly collinear columns,
a least-squares fit can reproduce almost any field through large, cancelling coefficients and so overstate how much
the basis explains.

\paragraph{Parametric map.} Each training field is projected onto $\mathbf{B}$ by least squares. Its uniform coefficient
$c_0$ enters a Gaussian process (GP) through the residual $c_0-m_a(\mathbf{p})$. The non-uniform part
$\mathbf{y}=\mathbf{Q}\mathbf{c}_Q$ (the rigid coefficients are discarded) is compressed by principal component
analysis over the training set, keeping the leading components that hold 99.9\% of its variance, at most eight; one GP
maps $\mathbf{p}$ to each score. The magnification carries most of the field energy, so it is modelled separately:
otherwise it would crowd the non-uniform modes out of the principal components. All GPs use a Mat\'ern-5/2 kernel with
one length scale per input plus a white-noise term, normalized targets, and hyperparameters fitted by maximum marginal
likelihood (scikit-learn). Inputs are the 16 parameters scaled to $[0,1]$ with the bounds of the parameter schema,
in logarithmic space for parameters swept on a log scale. The prediction for a new parameter set is
\begin{equation}\label{eq:basispred}
  \hat{\mathbf{f}}(\mathbf{p}) = \bigl(m_a(\mathbf{p})+\hat r_0(\mathbf{p})\bigr)\,\mathbf{U}_{00}
     + \bar{\mathbf{y}} + \sum_{k} \hat s_k(\mathbf{p})\,\mathbf{v}_k ,
\end{equation}
with $\hat r_0$ and $\hat s_k$ the GP means, $\bar{\mathbf{y}}$ the training mean and $\mathbf{v}_k$ the principal
directions. Every predicted field lies in the span of $\mathbf{B}$, so the reduced physics holds exactly (hard
constraint). We call this realization \emph{physics-structured}.

\subsection{Physics-informed realization: parametric PINN}\label{sec:pinn}
The second realization enforces the same reduced physics as a penalty instead of a basis. Three fully connected
networks with inputs $(\boldsymbol\xi,\mathbf{p})$ (four hidden layers of width 64, $\tanh$ activations) define
\begin{equation}
  \mathbf{u}_m = m_a\,\boldsymbol\xi + S\,N_u,\qquad e = m_a + S\,N_e,\qquad \phi = S\,N_\phi,\qquad
  \mathbf{f} = \mathbf{u}_m + \nabla\phi ,
\end{equation}
with a fixed output scale $S=\SI{30}{ppm}$. The analytic prior is the exact solution for $N_u=N_e=0$, so training starts
from it. The loss is the sum of four mean-squared terms, all divided by $S^2$:
\begin{itemize}
  \item \emph{data}: $\mathbf{f}$ against the reference field at the pads of the training runs;
  \item \emph{equilibrium}: the residual $\nabla\cdot\boldsymbol\sigma$ of~\eqref{eq:membrane} for $(\mathbf{u}_m,e)$ at
        random interior points;
  \item \emph{traction}: $\boldsymbol\sigma\mathbf{n}$ at random edge points;
  \item \emph{gauge}: the mean translation and mean rotation of the network part of $\mathbf{f}$ over the die, for
        random parameter sets, which fixes the rigid motion as in the reference data.
\end{itemize}
Derivatives with respect to $\boldsymbol\xi$ are taken by automatic differentiation (PyTorch). The collocation and edge
points are paired with parameter sets drawn from a pool that holds the training parameters and 4096 additional points
sampled uniformly in the unit cube, which have no reference run. This is the mechanism by which a PINN can add
information beyond the data: the physics terms constrain the model at parameter values where no simulation exists.
Training uses Adam with a learning rate of $2\times10^{-3}$, cosine-decayed to 1/50 of that value over 6000 steps, with
2048 data points, 1024 interior points and 256 edge points per step. As in Section~\ref{sec:reduced}, $e$ and $\phi$ are
free fields; the physics terms restrict how they combine into $\mathbf{f}$, not their values.

\subsection{Reconstruction from sparse overlay marks}\label{sec:marks}
In production the process parameters of a bond are known, but the overlay is measured only at a few marks per die. We
combine both in a Gaussian (maximum a posteriori) estimate within the span of the physics-structured model. Let
$\mathbf{M}=[\mathbf{R},\mathbf{U}_{00},\mathbf{v}_1,\dots,\mathbf{v}_r]$ collect the rigid modes, the uniform mode and the
model's principal directions, and write the field as
\begin{equation}
  \mathbf{f} = \hat{\mathbf{f}}(\mathbf{p}) + \mathbf{M}\mathbf{z},\qquad \mathbf{z}\sim\mathcal{N}(\mathbf{0},\boldsymbol\Sigma),
\end{equation}
with the prediction~\eqref{eq:basispred} as mean. The prior covariance $\boldsymbol\Sigma$ is diagonal: a broad
variance for the rigid modes, which the model does not predict, and for each other mode twice its five-fold
cross-validated prediction variance on the training runs. With $k$ marks measured with noise $\sigma$, $\mathbf{A}$ the
rows of $(L/2)\,\mathbf{M}$ at the marks and $\mathbf{y}_k$ the measured overlays in nm,
\begin{equation}\label{eq:map}
  \hat{\mathbf{z}} = \boldsymbol\Sigma\mathbf{A}^{\top}\bigl(\mathbf{A}\boldsymbol\Sigma\mathbf{A}^{\top}+\sigma^2\mathbf{I}\bigr)^{-1}
                     \bigl(\mathbf{y}_k-\tfrac{L}{2}\,\hat{\mathbf{f}}_k\bigr),\qquad
  \hat{\mathbf{d}} = \tfrac{L}{2}\bigl(\hat{\mathbf{f}}+\mathbf{M}\hat{\mathbf{z}}\bigr),
\end{equation}
where $\hat{\mathbf{f}}_k$ is the prior mean at the marks. The estimate needs no training beyond the parametric model;
for each die it is one small linear solve. With few marks it stays close to the prior; as marks are added it moves
toward the data, but only along directions the physics basis allows.

\section{Experiments}\label{sec:experiments}
All experiments train on subsets of the \NTrainRuns~training runs of Section~\ref{sec:reference} and test on the
\NHoldoutRuns~holdout runs, which no model sees during training. Unless stated otherwise, accuracy is the holdout $R^2$
of three scalar targets computed from each predicted field with the overlay fit~\eqref{eq:overlayfit}: the
magnification $m_x$, the non-uniform component $\mathrm{resid}_{3\sigma}$ and the largest misalignment $d_{\max}$.
Models that predict whole fields are also scored by the root-mean-square pad error over all holdout pads (field RMSE).

\paragraph{Models.} Besides the two realizations of Section~\ref{sec:method} we compare four baselines that use the
same training runs:
\begin{itemize}
  \item \emph{analytic only}: the field $m_a(\mathbf{p})\,\boldsymbol\xi$ from~\eqref{eq:prior}, which needs no
        simulation;
  \item \emph{GP, data only}: one GP per scalar target on the \NParams~scaled parameters, with no physics and no prior;
  \item \emph{analytic + GP residual} and \emph{analytic + MLP ensemble}: one GP per target, or an ensemble of
        multilayer perceptrons, trained on the residual between each scalar target and its closed-form estimate
        ($m_a$ for the magnifications, the corresponding $d_{\max}$ for the largest misalignment).
\end{itemize}
The two hybrid baselines thus receive the same analytic prior as our models; what they lack is the field structure.
They predict the scalar targets directly, so they have no field RMSE.

\paragraph{R0: representation.} Before any learning we ask how much of the reference overlay the reduced physics can
represent at all. For every training run we remove the affine part of the field (the industry six-parameter overlay
model) and project the remainder onto four candidate spans: the membrane responses of Section~\ref{sec:basis}, the
bending gradients, both together, and a generic polynomial field of the same total degree in each component. Each span
is orthonormalized together with the affine modes, so nearly collinear columns cannot inflate the fit. We report the
share of the non-affine energy each span captures and the median remaining RMS per die, as a function of polynomial
degree and hence dimension.

\paragraph{E1: learning curves.} Every model is trained on random subsets of $N$ training runs, $N$ = \EoneSizes, and
evaluated on the full holdout set. All models of one repeat share the same subset. We use \EoneRepsSmall~subsets per
size up to $N=\EoneFullRepsMaxN$, \EoneRepsFiveTwelve~at $N=\EoneMidN$ and \EoneRepsMax~at $N=\EoneMaxN$ (the last only for the models whose
training cost allows it), and report medians. The PINN is trained at $N$ = \PinnSizes, with \PinnReps~subsets each.

\paragraph{E2: reconstruction from sparse marks.} The physics-structured model is trained on \EtwoPriorSmall~or
\EtwoPriorLarge~runs (\EtwoReps~random subsets each) and serves as the prior of Section~\ref{sec:marks}. For every
holdout die we then simulate a measurement: the reference overlay at $k$ = \EtwoKs~mark positions (the die corners
and mid-edge or centre positions for the smaller counts, square grids for the larger ones) plus independent Gaussian
noise of standard deviation $\sigma$ = \EtwoNoises~nm on each component. From the $k$ marks and the known process
parameters each method reconstructs the full $\NPads\times\NPads$ field, which is compared with the reference field.
Besides the MAP estimate~\eqref{eq:map} and the prior alone (no marks), we evaluate four empirical corrections fitted
by least squares to the same marks: the six-parameter affine model used in production, second- and third-order
polynomial corrections in each component (where the number of marks allows), and the physics prior plus an affine fit
to its residual at the marks. The last one has the same prior information as our estimate and differs only in not
using the physics basis. We report the median field RMSE over the holdout dies.

\paragraph{E3: robustness.} Two departures from the clean setting of E1 are tested with $N=\EthreeN$ training runs
and \EthreeReps~subsets. First, independent Gaussian noise of \EthreeNoises~nm is added to every training field, as an
imperfect measurement would add it; the holdout set stays clean. Second, two extrapolation splits train only on runs
with $\Delta T_{\mathrm{head}}\le\SI{\EthreeDTThreshold}{\celsius}$ or with $L\le\SI{\EthreeLThreshold}{\milli\meter}$
and test only on the holdout runs beyond these limits (\EthreeOodDTNTest~and \EthreeOodLNTest~runs). The PINN is not
included in E3 because of its training cost.

\paragraph{E4: ablations.} Within the protocol of E1 we remove one ingredient at a time from the physics-structured
model: the analytic prior (the uniform coefficient is then learned from zero), or the bending columns (membrane
responses only). For the PINN we remove the physics loss terms, which leaves a network fitted to the data with the
same architecture and prior, or the bending potential; these ablations use $N=\PinnAblationN$ with
\PinnAblationReps~subsets.

\section{Results}
\begin{figure}[t]\centering\includegraphics[width=0.7\linewidth]{fig_r0_representation.pdf}
\caption{Share of the non-affine FEM overlay that each basis represents (degree 4).}\label{fig:r0}\end{figure}
\begin{figure*}[t]\centering\includegraphics[width=\linewidth]{fig_e1_learning_curves.pdf}
\caption{Holdout error vs number of FEM training runs (median and min--max over random subsets).}\label{fig:e1}\end{figure*}
\begin{figure*}[t]\centering\includegraphics[width=\linewidth]{fig_e2_sparse_marks.pdf}
\caption{Full-die reconstruction from $k$ overlay marks.}\label{fig:e2}\end{figure*}
\begin{figure}[t]\centering\includegraphics[width=\linewidth]{fig_e3_robustness.pdf}
\caption{Robustness to noisy training data and extrapolation.}\label{fig:e3}\end{figure}
\begin{figure}[t]\centering\includegraphics[width=\linewidth]{fig_e4_ablations.pdf}
\caption{Ablations.}\label{fig:e4}\end{figure}
\begin{figure*}[t]\centering\includegraphics[width=\linewidth]{fig_field_example.pdf}
\caption{Non-affine overlay of the holdout die with the largest non-uniform shift.}\label{fig:field}\end{figure*}
\begin{table}[t]\centering\small
\caption{Holdout accuracy versus number of FEM training runs (median over repeats).}\label{tab:e1}
\input{../../results/benchmark/figures/table_e1_summary.tex}
\end{table}

\subsection{Reconstruction from sparse overlay marks (E2)}\label{sec:res-e2}
Figure~\ref{fig:e2} separates two effects: what the physics prior contributes, and what the physics basis adds on
top of it. Numbers below are medians over the holdout dies at a measurement noise of \EtwoNoiseRef~nm.

\emph{With few marks, the prior does the work.} With $k=\EtwoKFour$ marks the six-parameter affine model leaves a field
error of \EtwoAffineFour~nm, because it cannot represent the non-affine part of the overlay. Adding the physics prior and
fitting the affine model only to its residual lowers the error to \EtwoPriorAffineFour~nm with a prior trained on
\EtwoPriorLarge~runs and \EtwoPriorAffineFourSmallPrior~nm with \EtwoPriorSmall~runs, a factor of
\EtwoPriorGainFourLarge~and \EtwoPriorGainFourSmall. Our MAP estimate gives no further gain at this mark count
(\EtwoOursFour~and \EtwoOursFourSmallPrior~nm): so few marks cannot resolve more than the prior already predicts. Without
any marks the \EtwoPriorLarge-run prior alone leaves \EtwoPriorOnly~nm, the same as the affine model with
\EtwoKFour~marks.

\emph{With dense marks, the physics basis adds value.} With $k=\EtwoKSixteen$ marks the MAP estimate reaches
\EtwoOursSixteen~nm (\EtwoPriorLarge-run prior) and \EtwoOursSixteenSmallPrior~nm (\EtwoPriorSmall-run prior), against
\EtwoPriorAffineSixteen~and \EtwoPriorAffineSixteenSmallPrior~nm for the prior with an affine correction: the physics
basis lowers the error by a further factor of \EtwoBasisGainSixteenLarge~to \EtwoBasisGainSixteenSmall. The
third-order correction, the strongest empirical model these marks support, reaches \EtwoPolyThreeSixteen~nm without any
prior. At $k=\EtwoKNine$ the benefit depends on the prior: it is a factor of \EtwoBasisGainNineSmall~with the
\EtwoPriorSmall-run prior, but with the \EtwoPriorLarge-run prior the MAP estimate (\EtwoOursNine~nm) is no better than
the prior with an affine correction (\EtwoPriorAffineNine~nm), and at \EtwoNoiseHigh~nm noise it is slightly worse
(\EtwoOursNineNoiseHigh~vs \EtwoPriorAffineNineNoiseHigh~nm).

For a fab this splits into two recommendations: a prior from a few hundred simulations is worth having even with the
minimum number of marks, and the physics basis becomes worth its added complexity once dense marking is available.

\section{Discussion and limitations}\label{sec:discussion}
\begin{itemize}
  \item No experimental overlay data: validation is against an independent, higher-fidelity model (3D transient FEM with
        bond-wave locking) whose physics is richer than the reduced model, so the test is not an inverse crime; the
        reduced-physics gap is quantified (R0).
  \item Plausibility: magnitudes consistent with published D2W measurements (die-corner scaling 80--150~nm \cite{bex2024}).
  \item What a fab would measure to validate: \todo{protocol --- marks per die, dies per wafer, recipe split}.
  \item Hard constraint vs PINN: with the eigenstrain left free (closure field), the membrane residual only
        restricts the \emph{shape} of the field, which the curl-free bending part can absorb anyway; the
        information that transfers to unseen process points is the analytic prior and the low-dimensional span,
        which the hard-constrained model encodes exactly and a GP exploits efficiently in 16 dimensions.
        PINN ablation at $N=64$: field RMSE 92~nm with the physics loss, 88~nm without (Fig.~\ref{fig:e4}).
        A PINN becomes worthwhile when the loads are known in closed form (e.g. a solved thermal field), not as
        a closure model.
\end{itemize}

\section*{Code and data availability}
Repository, study YAML files, pinned DOLFINx image digest, and every benchmark CSV. \todo{URL}

\bibliographystyle{unsrt}
\bibliography{refs}
\end{document}
